\chapter{Reasoning, search, and planning}
\section{Planning as choosing actions under dependencies}
The equipment request becomes more difficult when a useful result requires several dependent choices. A camera must be compatible with a power supply, available for the trip dates, and allowed under the student's training status. Planning organizes possible actions so the system can reach a valid end state without unnecessary work.

A plan is a proposed structure for future action. It is not evidence that the actions occurred. A useful plan contains preconditions, expected results, and dependencies. “Check eligibility, then prepare a reservation” is stronger than “handle the booking” because the second step depends on the first. The runtime still has to verify eligibility when the plan executes.

Reactive control chooses a next step from the latest observation. Plan-and-execute constructs a broader sequence before beginning, then revises it when observations invalidate assumptions. Neither dominates every task. A short lookup may not justify a planning phase; a multi-resource expedition may benefit from making dependencies explicit.

\section{Reasoning paths and self-consistency}
A model can generate different intermediate paths to an answer. Self-consistency samples multiple reasoning paths and aggregates their answers; Wang and colleagues report improvements on selected reasoning benchmarks \cite{self-consistency}. The aggregation rule requires comparable final answers. Voting over free-form essays without a normalization rule can confuse wording differences with substantive disagreement.

Why might voting help? Under a deliberately restrictive model, suppose three independent attempts each produce the correct binary answer with probability \(p\). Majority correctness occurs when exactly two or all three are correct:

\begin{equation}
P(\text{majority correct})=3p^2(1-p)+p^3.
\end{equation}

At \(p=0.7\), this equals 0.784. The improvement follows from the assumed independence and binary outcome structure. Real model outputs can share errors because they use the same model, prompt, and evidence. If every attempt follows the same wrong policy passage, repeated sampling does not create independent evidence.

The correct experimental question is therefore whether voting improves this task under a declared resource budget. Record individual outputs, normalize answers using a documented rule, and include ties or invalid answers in the procedure. Report how much additional inference the voting design consumes.

\section{Search over partial solutions}
Tree of Thoughts explores candidate intermediate states using generation and evaluation within a search process \cite{tree-thoughts}. A general search algorithm needs a state representation, successor function, goal test, and selection rule. For equipment planning, a state could be a partial kit with unresolved requirements. A successor adds a compatible item. The goal test checks that every requirement is satisfied and all constraints hold.

Breadth-first search explores states in order of depth. Depth-first search follows one branch further before backtracking. Beam search retains a limited number of candidates at each stage according to a score. These strategies differ in memory use and which alternatives they discard. A model-based evaluator can rank candidates, but an incorrect score may prune the only valid solution.

If a full tree has branching factor \(b>1\) and depth \(d\), its number of nodes is:

\begin{equation}
1+b+b^2+\cdots+b^d=\frac{b^{d+1}-1}{b-1}.
\end{equation}

For \(b=3\) and \(d=4\), there are 121 nodes. This counts a complete tree under fixed branching, not an actual model workload. Pruning and early stopping can reduce the explored set, while retries and evaluation calls add work not visible in the node count.

\section{Worked example: a constrained kit search}
Suppose a kit needs one camera and one battery, with a total weight limit of 3 units. Camera A weighs 2 units and uses battery X, which weighs 2. Camera B weighs 1 unit and uses battery Y, which weighs 1. The first partial state chooses A because it has the highest quality score. Extending it with X creates a 4-unit kit, which violates the hard limit.

A greedy strategy that refuses to revise the camera choice gets stuck. A search strategy can backtrack and select B plus Y, producing a valid 2-unit kit. The quality score is useful only among feasible plans. A hard constraint should not be overridden because an attractive partial state received a high model score.

The example also shows why evaluating partial states is difficult. A looks better before compatibility and weight are considered together. A good evaluator should account for whether a partial state can still be completed, not only how appealing its current contents look. In a real application, exact compatibility and weight checks can often be deterministic, leaving the model to interpret preferences or propose candidates.

\section{Reflection and refinement}
Reflection uses feedback from an attempt to guide a later attempt. Reflexion studies agents that retain verbal feedback, while Self-Refine studies iterative generation, feedback, and refinement \cite{reflexion,self-refine}. These methods differ from searching a broad tree: they revise a trajectory or candidate through feedback rather than necessarily keeping many alternatives alive.

Feedback quality determines what the loop can learn. “The kit is poor” is vague. “The selected battery is incompatible with camera A under catalog record C12” identifies a correctable constraint. A deterministic test failure can anchor revision. A self-generated critique without new evidence may merely change style or introduce an unsupported alternative.

Store the reason for a revision, the changed fields, and the test result after revision. If the model changes unrelated parts of the plan, rerun the relevant tests. Avoid declaring progress based on a more confident explanation. The acceptance condition is a property of the plan or outcome.

\section{Test-time compute and reasoning models}
Training changes model parameters; test-time computation spends resources on a particular request. Sampling multiple candidates, evaluating branches, or revising an answer are application-level ways to spend test-time compute. Some model families also perform extended internal reasoning. These are different mechanisms, even when both increase request cost or latency.

The application designer needs an allocation policy: which tasks justify extra work, how the extra work is bounded, and how benefit is measured. A routing rule might use a direct method for exact lookup and a bounded search for a multi-constraint plan. But a difficulty estimate can be wrong, so the fallback and stop behavior must also be evaluated.

A model's verbal explanation of its reasoning should not be treated as an execution trace of internal computation. Counterfactual experiments can test whether a specified input edit changes output behavior \cite{chive}. They do not reveal every internal cause. In the course, require concise justifications tied to evidence and tests rather than reward the length of generated reasoning.

\section{Planning with observations}
An initial plan may assume that C18 is available. After an inventory read contradicts that assumption, the plan must change. Replanning should preserve completed valid work and invalidate dependent steps. If only the camera choice changes, a policy check about the user's training may remain valid, while compatibility checks for the old camera must be repeated.

Represent dependencies explicitly so that invalidation is systematic. A directed acyclic graph is useful for a fixed plan with no cycles. A reactive workflow may include loops for repair, but those loops need counters and stopping conditions. Graph structure makes dependencies inspectable; it does not ensure that the node descriptions are correct.

A plan should also distinguish reversible preparation from effects. Searching, calculating, and drafting can usually precede approval. Commitment requires current validation and authorized execution. An elegant plan is still incomplete until the system can explain which effects occurred and which remain proposals.

\section{Exercises}
\begin{enumerate}
\item Compute three-vote majority correctness at \(p=0.6\) under the chapter's independent binary model. Explain why this is not a prediction for three real model calls.
\item Count the nodes in a full tree with branching factor 2 and depth 5. State what extra model work the count omits.
\item Extend the kit example with a third camera and battery. Construct a case where greedy quality ranking fails but backtracking finds a feasible kit.
\item Write a reflection message grounded in a failed compatibility test. Specify which fields may change and which tests must rerun.
\item Compare direct answering, three-sample voting, and one revision on a fixed task set. Predeclare the budget and how ties, invalid responses, and timeouts are scored.
\end{enumerate}

\section{Further study and laboratory connection}
Read \cite{self-consistency,tree-thoughts,reflexion} as distinct ways to organize additional computation. Lab 7 provides the experimental structure for paired comparisons. An acceptable conclusion is that more computation failed to improve the outcome, provided the comparison and accounting are explicit.

